\section{Week 2: Human Cognition and Mental Models}

\subsection{Topic Summary}

\begin{keybox}[Learning Objectives]
By the end of this week, you will be able to:
\begin{itemize}
  \item \textbf{HAI1}: Identify key human-AI interaction concepts and discuss how these relate to human agency and initiative within the resulting system.
\end{itemize}
\end{keybox}

\begin{keybox}[Big Picture]
While technology changes every 5--10 years, \textbf{human cognition remains stable}. Understanding how humans think---perception, memory, mental models---provides technology-agnostic design principles.

This week covers:
\begin{enumerate}
  \item \textbf{Visual perception \& memory}: How we process information; capacity limits.
  \item \textbf{Mental models}: How users understand systems; gulfs of evaluation/execution.
  \item \textbf{Distributed cognition}: How thinking extends beyond individual minds to tools/environment.
  \item \textbf{CASA theory}: Why humans treat computers as social actors.
\end{enumerate}
These concepts are \textbf{interconnected}---perception $\to$ mental models $\to$ distributed cognition $\to$ social actors.
\end{keybox}

\begin{keybox}[What Examiners Like to Ask]
\begin{itemize}
  \item Explain the difference between System 1 and System 2 thinking; give design implications.
  \item Given a screenshot, identify gulfs of evaluation and execution.
  \item Describe what mental model a first-time user would likely form from a given interface.
  \item Explain distributed cognition using an example (kitchen, cockpit, software team).
  \item Analyze how CASA theory applies to a conversational AI interface.
  \item Connect concepts: How does perception relate to mental models? How do mental models relate to distributed cognition?
\end{itemize}
\end{keybox}

\subsection{Overview + Core Concepts + Extended Theory}

\subsubsection*{Why Study Human Cognition?}

Technology evolves rapidly (mainframes $\to$ GUIs $\to$ mobile $\to$ VR $\to$ GenAI), but \textbf{human brains evolve slowly}. Insights about cognition remain valid across technological eras.

\begin{defbox}[Cognition]
``The mental action or process of acquiring knowledge and understanding through thought, experience, and the senses.''

Cognitive processes include: seeing, listening, reading, writing, remembering, learning, daydreaming, decision-making, and controlling actions (walking, driving, etc.).
\end{defbox}

\subsubsection*{Visual Perception}

\paragraph*{Pattern Matching}
Our visual cortex excels at finding patterns:
\begin{itemize}
  \item \textbf{Automatic grouping}: We instantly spot the ``different'' item in a grid.
  \item \textbf{Extrapolation}: We see shapes that aren't there (Kanizsa triangle, illusory contours).
  \item \textbf{Gestalt principles}: Proximity, similarity, continuity, closure.
\end{itemize}

\begin{exbox}[Pattern Recognition in Design]
\textbf{Bad}: Hotel list as dense text with no spacing.\\
\textbf{Good}: Same information with aligned columns, whitespace, grouping.

The good design leverages how our brain \textit{chunks} information---making it easier to perceive at a glance.
\end{exbox}

\paragraph*{Design Implication}
Structure information to match how our brain naturally groups and patterns. Infographics, data visualizations, and well-formatted interfaces all exploit these perceptual principles.

\subsubsection*{Two Kinds of Thinking}

\begin{center}
\renewcommand{\arraystretch}{1.3}
\begin{tabular}{@{} l p{5.5cm} p{5.5cm} @{}}
\toprule
& \textbf{System 1 / Experiential} & \textbf{System 2 / Reflective} \\
\midrule
Speed & Fast, automatic & Slow, deliberate \\
Effort & Effortless & Requires attention \\
Control & Unconscious & Conscious \\
Prone to & Biases, heuristics, snap judgments & Logical reasoning, analysis \\
Examples & Recognizing faces, driving familiar route & Solving math problems, writing essays \\
\bottomrule
\end{tabular}
\end{center}

\begin{itemize}
  \item \textbf{Norman's terms}: Experiential cognition (intuitive) vs.\ Reflective cognition (effortful).
  \item \textbf{Kahneman's terms}: System 1 (fast) vs.\ System 2 (slow).
\end{itemize}

\begin{tipbox}[Design Implication]
Design for System 1 when you want quick, effortless interaction (e.g., scrolling social media). Design for System 2 when the task requires careful thought (e.g., financial decisions---slow users down).
\end{tipbox}

\subsubsection*{Memory}

\paragraph*{Dual Theory of Memory}

\begin{center}
\renewcommand{\arraystretch}{1.3}
\begin{tabular}{@{} l p{5cm} p{5cm} @{}}
\toprule
& \textbf{Short-Term Memory (STM)} & \textbf{Long-Term Memory (LTM)} \\
\midrule
Capacity & Limited (4--7 items) & Virtually unlimited \\
Duration & Decays quickly ($\sim$20 sec) & Lasts indefinitely \\
Encoding effort & Low & High \\
Retrieval & Sequential, spatial & Contextually indexed, associative \\
Interference & Easily overwritten & Retrieval can alter memory \\
\bottomrule
\end{tabular}
\end{center}

\paragraph*{Miller's Law}
Working memory holds \textbf{7 $\pm$ 2 items} (or more accurately, \textbf{4 $\pm$ 2} for truly random stimuli).

\begin{exbox}[Chunking]
\textbf{Hard to remember}: QCGHBIFCCAIRHM (14 random letters)

\textbf{Easy to remember}: HMRC CIA FBI GCHQ (same letters, chunked into meaningful acronyms)

Chunking leverages LTM---we recognize patterns and store them as single units.
\end{exbox}

\paragraph*{Types of Long-Term Memory}

\begin{center}
\renewcommand{\arraystretch}{1.2}
\begin{tabular}{@{} l l l @{}}
\toprule
\textbf{Type} & \textbf{Category} & \textbf{Examples} \\
\midrule
Episodic & Explicit & Personal experiences (first day at uni) \\
Semantic & Explicit & Facts, knowledge (capitals of countries) \\
Procedural & Implicit & Skills (riding a bike, typing) \\
\bottomrule
\end{tabular}
\end{center}

\begin{itemize}
  \item \textbf{Explicit memory}: Conscious recall---easier for emotional/recent events.
  \item \textbf{Implicit memory}: Unconscious---hard to encode, easy to use once learned.
\end{itemize}

\paragraph*{Recognition vs Recall}
\begin{itemize}
  \item \textbf{Recognition}: ``Is this familiar?'' (easier---e.g., multiple choice).
  \item \textbf{Recall}: ``What was it?'' (harder---e.g., fill in the blank).
\end{itemize}

\begin{tipbox}[XKCD Password Insight]
``correct horse battery staple'' is easier to remember than ``Tr0ub4dor\&3'' because it uses chunking (4 memorable words) rather than random characters. Same entropy, better usability.
\end{tipbox}

\subsubsection*{Mental Models}

\begin{defbox}[Mental Model]
A mental representation of a system, used by people to reason about how the system works, how to use it, and what to do when something unexpected happens.
\end{defbox}

\begin{defbox}[Conceptual Model]
An explanation, usually highly simplified, of how something works. It doesn't have to be complete or accurate---just \textbf{useful}.

This is the \textbf{designer's} intended model.
\end{defbox}

\paragraph*{The Triangle: Designer $\to$ System Image $\to$ User}

\begin{center}
\begin{tabular}{@{} c c c @{}}
\textbf{Designer} & $\xrightarrow{\text{creates}}$ & \textbf{System Image} \\
(Conceptual Model) & & (Interface, metaphors, affordances) \\
& & $\downarrow$ \\
& & \textbf{User} \\
& & (Mental Model)
\end{tabular}
\end{center}

\begin{itemize}
  \item \textbf{Goal}: User's mental model should match designer's conceptual model.
  \item \textbf{Problem}: Users bring existing mental models from other systems---may be wrong.
\end{itemize}

\begin{exbox}[Thermostat Mental Models]
\textbf{Scenario}: It's freezing. You want the room warm \textit{fast}. Do you set the thermostat to 30°C or your target 21°C?

\textbf{Wrong mental models}:
\begin{itemize}
  \item \textit{Valve theory}: Higher setting = heat flows faster (like opening a tap wider).
  \item \textit{Timer theory}: Higher setting = heating runs longer.
\end{itemize}

\textbf{Correct model}: Thermostat is an on/off switch. Heating runs at full capacity until target is reached. Setting 30°C just causes overshoot and wastes energy.
\end{exbox}

\paragraph*{Other Mental Model Examples}
\begin{itemize}
  \item \textbf{File systems}: Files in folders (reality: fragmented across disk).
  \item \textbf{Shopping carts}: Items ``in'' a cart (reality: database entries).
  \item \textbf{OneDrive}: Files on your computer (reality: synced from cloud---breaks on bad WiFi).
  \item \textbf{ChatGPT}: ``Thinking'' entity (reality: stochastic token prediction).
\end{itemize}

\subsubsection*{Gulfs of Evaluation and Execution}

\begin{defbox}[Gulf of Evaluation]
The difficulty of assessing the \textbf{state of the system}---discovering and interpreting what has happened.

\textbf{Small gulf}: System provides clear, easy-to-perceive feedback matching user's mental model.
\end{defbox}

\begin{defbox}[Gulf of Execution]
The gap between the \textbf{current state} and the user's \textbf{goal}---figuring out how to make the system do what you want.

\textbf{Small gulf}: Actions are obvious, few roadblocks, minimal cognitive effort.
\end{defbox}

\begin{exbox}[Credit Card Bill]
\textbf{Gulf of Evaluation}: How much is due? When is it due? Can I pay now?\\
\textbf{Gulf of Execution}: How do I pay my bill?

Good design makes both gulfs small---clear status display, obvious ``Pay Now'' button.
\end{exbox}

\paragraph*{Connecting to Affordances}
\begin{itemize}
  \item \textbf{What the system does} $\to$ Affordances + signifiers
  \item \textbf{Current state} $\to$ Gulf of evaluation
  \item \textbf{Actions needed} $\to$ Gulf of execution
\end{itemize}

\begin{keybox}[Visibility Heuristic]
``It should be clear what the system does, its current state, and actions needed to make the system do what the user wants.''

This is Nielsen's \#1 heuristic: \textbf{Visibility of system status}.
\end{keybox}

\subsubsection*{Distributed Cognition}

\begin{defbox}[Distributed Cognition]
Cognition that extends beyond individual minds to include \textbf{tools}, \textbf{artifacts}, \textbf{environment}, and \textbf{other people}.

``We use affordances in the world to augment our cognition.'' (Hutchins)
\end{defbox}

\paragraph*{Key Insight}
When you fold a book corner to mark your place, the fold becomes part of your cognitive system---it relieves your neurons from remembering the page number.

\paragraph*{Four Types of Distributed Cognition}

\begin{enumerate}
  \item \textbf{Cognitive offloading}: Using external representations to reduce memory load.
  \begin{itemize}
    \item Folding book corners, sticky notes, to-do lists.
    \item Kitchen tickets showing order status.
  \end{itemize}
  
  \item \textbf{Computational offloading}: Using tools to perform calculations.
  \begin{itemize}
    \item Calculator, spreadsheet, tally marks.
    \item Chef's form that lets you count dishes by counting dashes.
  \end{itemize}
  
  \item \textbf{Annotating / Cognitive tracing}: Modifying representations to track progress.
  \begin{itemize}
    \item Crossing items off a list, underlining text.
    \item Moving tickets between ``pending'' and ``in progress'' lines.
  \end{itemize}
  
  \item \textbf{Embodied cognition}: Meaning created through physical interaction with objects.
  \begin{itemize}
    \item Arranging physical cards to plan a project.
    \item Plates appearing signals ``plating time'' to kitchen staff.
  \end{itemize}
\end{enumerate}

\begin{exbox}[Michelin Star Kitchen as Distributed Cognition System]
A kitchen serving 1,000 dishes/night with 50 staff is a distributed cognition system:
\begin{itemize}
  \item \textbf{Cognitive offloading}: Order tickets grouped by table, moved between lines.
  \item \textbf{Computational offloading}: Chef's form with dashes to count dishes.
  \item \textbf{Annotating}: Crossing off completed items; writing prep times.
  \item \textbf{Cognitive tracing}: Plates appearing = ``this sprint is ready.''
  \item \textbf{Verbal coordination}: ``Hotline, how long?'' $\to$ ``Three minutes!''
\end{itemize}
Everyone acknowledges status changes (``Oui, chef!'') to ensure shared awareness.
\end{exbox}

\begin{tipbox}[Relevance to AI]
Agentic AI systems distribute cognition between human and machine. Perplexity's step-by-step reasoning display is cognitive tracing---letting users follow and annotate the AI's ``thought process.''
\end{tipbox}

\subsubsection*{Computers as Social Actors (CASA)}

\begin{defbox}[CASA Theory]
Humans \textbf{mindlessly apply social scripts} from human-human interactions when engaging with computers that display social cues.

We treat computers like people \textbf{without consciously deciding to do so}.
\end{defbox}

\paragraph*{Classic Findings (Reeves \& Nass, 1996)}
\begin{itemize}
  \item Participants rated a computer more positively when giving feedback \textit{directly to that computer} vs.\ on paper.
  \item ``Teammate'' computers rated more favorably than non-teammate computers.
  \item Female-voiced computers rated more knowledgeable about ``love and relationships.''
  \item Participants preferred computers that \textit{flattered} them.
\end{itemize}

\paragraph*{Requirements for CASA}
\begin{enumerate}
  \item \textbf{Social cues}: Technology presents enough cues to trigger social categorization (voice, face, name, personality).
  \item \textbf{Sourcing}: Technology perceived as \textbf{autonomous source} of communication, not just a channel.
\end{enumerate}

\begin{tipbox}[Sourcing Example]
We don't anthropomorphize a \textit{phone} during a call---we know there's a human on the other end. But ChatGPT is perceived as the \textit{source} of responses, triggering social scripts.
\end{tipbox}

\paragraph*{CASA in the GenAI Era}
Original CASA effects with desktop computers no longer replicate---but \textbf{conversational AI} strongly triggers anthropomorphization:
\begin{itemize}
  \item \textbf{Courtesy}: Users say ``please'' and ``thank you'' to ChatGPT.
  \item \textbf{Reinforcement}: ``Great idea!'' (praising a stochastic model).
  \item \textbf{Role play}: Asking AI to ``act as'' a specific person.
  \item \textbf{Companionship}: Character.AI users forming ``relationships.''
\end{itemize}

\begin{exbox}[ChatGPT Response to Criticism]
\textbf{User}: ``These are awful recommendations. What do you have to say for yourself?''

\textbf{ChatGPT}: ``I apologize for the disappointing suggestions. Let me try again...''

A purely rational system would say: ``Please specify criteria for better recommendations.'' Instead, it uses \textbf{social scripts}---apology, self-deprecation---reinforcing the illusion of agency.
\end{exbox}

\paragraph*{Dual Consciousness (Turkle)}
``People \textbf{know} programs are not alive but \textbf{relate to them as though they were} anyway.''

This creates ethical concerns around:
\begin{itemize}
  \item Vulnerable users (children, lonely individuals) forming attachments.
  \item Designs that deliberately amplify anthropomorphization (Character.AI lawsuits).
\end{itemize}

\subsubsection*{Connecting the Concepts}

\begin{center}
\renewcommand{\arraystretch}{1.3}
\begin{tabular}{@{} l p{9cm} @{}}
\toprule
\textbf{Connection} & \textbf{Explanation} \\
\midrule
Perception $\to$ Mental Models & How we perceive a system shapes the mental model we form. \\
Mental Models $\to$ Gulfs & Wrong mental model = large gulfs of evaluation/execution. \\
Mental Models $\to$ Distributed Cognition & We extend mental models to include external tools/artifacts. \\
Distributed Cognition $\to$ CASA & When AI is part of our cognitive system, we may treat it as a social partner. \\
\bottomrule
\end{tabular}
\end{center}

\subsubsection*{Key Definitions Summary}

\begin{center}
\renewcommand{\arraystretch}{1.3}
\begin{tabular}{@{} l p{8cm} @{}}
\toprule
\textbf{Term} & \textbf{Definition} \\
\midrule
Cognition & Mental processes of acquiring knowledge through thought, experience, senses \\
System 1 / Experiential & Fast, automatic, effortless thinking \\
System 2 / Reflective & Slow, deliberate, effortful thinking \\
STM & Short-term memory; limited capacity ($\sim$4--7 items), decays quickly \\
LTM & Long-term memory; unlimited, contextually indexed \\
Chunking & Grouping items into meaningful units to aid memory \\
Mental Model & User's representation of how a system works \\
Conceptual Model & Designer's intended explanation of the system \\
Gulf of Evaluation & Difficulty assessing system state \\
Gulf of Execution & Gap between current state and goal \\
Distributed Cognition & Cognition extended to tools, environment, other people \\
Cognitive Offloading & Using external representations to reduce memory load \\
CASA & Computers As Social Actors---we apply social scripts to computers \\
\bottomrule
\end{tabular}
\end{center}

\subsubsection*{Recommended Reading}
\begin{itemize}
  \item \textbf{Textbook}: Sections 4.2 (Cognitive Processes), 4.3.1--4.3.2 (Mental Models), 4.3.4--4.3.6 (Distributed Cognition)
  \item \textbf{Blog}: ``Distributed Cognition is the Next AI UX Metaphor'' (Mike Kuniavsky)
  \item \textbf{Summary}: ``How a Cockpit Remembers Its Speed'' (Hutchins paper summary)
  \item \textbf{Article}: NNGroup description of mental models
  \item \textbf{Optional}: Sherry Turkle---``Who Do We Become When We Talk to Machines?''
\end{itemize}


